两个 0--1 变量的联合分布可以用一张 \(2\times2\) 表格写出：

\[
\begin{array}{c|cc|c}
& Y=0 & Y=1 & \\
\hline
X=0 & 0.4 & 0.1 & 0.5 \\
X=1 & 0.2 & 0.3 & 0.5 \\
\hline
& 0.6 & 0.4 & 1
\end{array}
\]

从联合表可以读出：边缘分布 \(P(X=1)=0.5\)，条件分布 \(P(Y=1\mid X=1)=0.3/0.5=0.6\)，协方差 \(\Cov(X,Y)=0.3-0.5\times0.4=0.1>0\)。这就是联合分布的价值——单看边缘，只能知道 \(Y\) 有 40\% 概率为 1；联合表才揭示 \(Y\) 为 1 的概率在 \(X=0\) 时是 20\%，在 \(X=1\) 时是 60\%。

知道身高分布和体重分布，并不能回答``高个子是否通常更重''；知道两项资产各自的波动，也不能计算组合风险。边缘分布只描述每个变量单独怎样取值，联合分布才保留它们如何配对出现。

本单元从二维表格和密度区域开始，逐步进入组合与几何：

\begin{enumerate}
\tightlist
\item
  联合、边缘与条件分布：从联合概率恢复单变量分布，并在观察一个变量后切片。
\item
  随机变量的独立与依赖：用联合分布能否分解来判断独立，而不是只看相关系数。
\item
  随机变量之和与卷积：理解同一个和可以由哪些取值对产生。
\item
  随机向量与协方差矩阵：把方差推广成方向相关的波动几何。
\item
  多维变量变换与 Jacobian：坐标变换怎样压缩或拉伸概率密度。
\end{enumerate}

\subsection{一个具体算例}

上面的 \(2\times2\) 表格中：边缘满足 \(P(X=1)P(Y=1)=0.5\times0.4=0.2\)，但 \(P(X=1,Y=1)=0.3\ne0.2\)，故 \(X,Y\) 不独立。

独立情形：\(X,Y\sim\iid\Bernoulli(0.5)\)，则 \(P(X=0,Y=0)=0.25\)，\(P(X=0,Y=1)=0.25\)，\(P(X=1,Y=0)=0.25\)，\(P(X=1,Y=1)=0.25\)。此时 \(X+Y\) 的分布是 \(\Binomial(2,0.5)\)——两个独立 Bernoulli 之和的卷积。

连续情形：\(X,Y\sim\iid\Uniform(0,1)\)，独立。\(Z=X+Y\) 的密度是三角分布：\(f_Z(z)=z\) 对 \(z\in[0,1]\)，\(f_Z(z)=2-z\) 对 \(z\in[1,2]\)。

\subsection{怎么读这一章}

首次阅读先抓前三篇：联合分布保存配对信息，条件分布是在观测后重新切片，独立性是整个联合分布的分解条件；卷积只在独立等条件下由边缘分布单独决定。后两篇把同一结构写成线性代数与变量替换：协方差只描述二阶依赖，Jacobian 负责概率质量在坐标变化下守恒。多分支变换和退化协方差可第二遍处理，但``不相关不等于独立''与``边缘不能恢复联合''不能后置。

\subsection{本章篇目}

以下篇目按概念依赖排列；若从具体问题进入，也可以先打开最接近当前困惑的一篇，再沿篇内链接回补。

\begin{itemize}
\tightlist
\item \hyperref[sec:05-Probability/05-Joint-Distributions/01]{``联合、边缘与条件分布''}；
\item \hyperref[sec:05-Probability/05-Joint-Distributions/02]{``随机变量的独立与依赖''}；
\item \hyperref[sec:05-Probability/05-Joint-Distributions/03]{``随机变量之和与卷积''}；
\item \hyperref[sec:05-Probability/05-Joint-Distributions/04]{``随机向量与协方差矩阵''}；
\item \hyperref[sec:05-Probability/05-Joint-Distributions/05]{``多维变量变换与 Jacobian''}。
\end{itemize}
